\documentclass{article}
\usepackage[utf8]{inputenc}
\usepackage[T1]{fontenc}
\usepackage{geometry}
\geometry{
    a4paper,
    left=2.5cm,
    right=2.5cm,
    top=2.5cm,
    bottom=2.5cm
}
\usepackage{graphicx}
\usepackage{float}
\usepackage{hyperref}
\usepackage[english]{babel}
\usepackage{amsmath}
\usepackage{amssymb}
\usepackage{array}
\usepackage{booktabs}
\usepackage{caption}
\usepackage{subcaption}
\usepackage{enumitem}
\usepackage{fancyhdr}
\usepackage{titlesec}
\usepackage{url}
\usepackage{textcomp}
\usepackage{xcolor}
\usepackage{listings}
\usepackage{multirow}
\usepackage{longtable}
\usepackage{tabularx}
\usepackage{times}

% Set up page style
\pagestyle{fancy}
\fancyhf{}
\rfoot{\thepage}
\renewcommand{\headrulewidth}{0pt}

% Title formatting
\titleformat{\chapter}[display]{\normalfont\bfseries}{\chaptertitle\thechapter}{0pt}{\Huge}
\titlespacing*{\chapter}{0pt}{50pt}{40pt}

% Document info
\title{Template Document}
\author{Academic Support Agent}
\date{\today}

\begin{document}
\maketitle

%\pagecolor{yellow!5}
%%%%%%% Print marks in left margin %%%%%%%
\pointsinmargin % To print points in margin
\boxedpoints   % or \bracketedpoints

%%%%%%% Include the exam cover defined by the file eaxm_cover.tex %%%%%%%
% Content from exam_cover would be included here


\newpage

%%%%%%% Begin exam questions %%%%%%%
\begin{questions}
%%Question 1
\question \textbf{Answer the following} 
    \begin{parts}
        \item Choose the correct answer of the following
        \begin{parts}
            \part[1] If $x = 2$ and $y = 5 $, then $x+y$ \\
            \begin{oneparchoices}
                \choice 2
                \CorrectChoice 7
                \choice 21
                \choice 1
            \end{oneparchoices}
            
            \part[1] If $x = 12$ and $y = 5 $, then $x+y$ \\
            \begin{oneparchoices}
                \CorrectChoice 17 \\
                \choice $\displaystyle 7$ \\ 
                \choice $\displaystyle 21$ \\
                \choice $\displaystyle 1$ 
            \end{oneparchoices}
            
            \part[2] Who invented the pencil
            \begin{oneparcheckboxes}
                \choice Me
                \choice You
                \choice They
                \choice Him
                \CorrectChoice Socrates
            \end{oneparcheckboxes}
        \end{parts}
    
        \part[2] If $f(x) = \sin(x)$, then $f'(x) = $\fillin.
        \part True or False
            \begin{parts} 
                \part[2] \tf[T] The world is all that is the case.
                \part[1] \tf[F] My favorite color is blue.
            \end{parts}
    \end{parts}
%

\singlespacing

%%Question 2
\question \textbf{Let $\displaystyle f(x) = \sin(x) + x^2$}
    \begin{parts}
        \part[4] Compute $\displaystyle \frac{df}{dx} $.\\
            \fillwithdottedlines{1.5in}
              %%%% To add solutions, add the following and uncomment the seconed line
              %%%% in the document \documentclass[12pt,addpoints,answers]{exam} %     To print with answers
            \begin{solution}
                $f'(x) = \cos(x) + 2 x $
            \end{solution}
        \part[5] Compute $\displaystyle\int_0^1 f(x) \, dx$.\\
            \fillwithdottedlines{1.5in}
            \begin{solution}
                Put here the answer...
            \end{solution}
    \end{parts}

%%Question 3
\question[10] \textbf{Describe the effect of error propagation on numerical results.}\\
\fillwithdottedlines{2.8in}


%%Question 4
\question
    \begin{parts}
        \part[3] What do you do with $f(x) = x$?
        \fillwithdottedlines{2.8in}
        
        \part[3] Is your answer different if $f(x) = \tan(x)$?
        \fillwithdottedlines{1.9in}
    \end{parts}

%%Question 4
\question[8]
    In no more than one paragraph, explain why the earth is round.
    \makeemptybox{1in}

%%Question 4
\question Answer the following 
    \begin{parts}
        \part[1] \label{eqn1}
            find $f(2.25)$.
            \answerline
        
        \part[1]\label{eqn2}
            Approximate $f'(0)$.
            \answerline
        
        \part[1]\label{eqn3}
            Approximate $\displaystyle \int_{1}^{5} f(x) \: \text{dx}$.
            \answerline

        \part[5] Redo questions \eqref{eqn1}, \eqref{eqn2} and \eqref{eqn3} with $f(1) = 4$ and $f(2) = 5$.
        \fillwithdottedlines{1.2in}
    \end{parts}

\end{questions}

\begin{center}
\fbox{\fbox{\parbox{5.5in}{\centering
Hope you all the best! \\
Dr. X Y Z, \quad \quad
Dr. M N T}}}
\includegraphics[width=0.25\textwidth]{logo}
\end{center}
\end{document}